\section{Notation}

\subsection{Metavariables}

\begin{longtable}[]{@{}rl@{}}
\toprule\noalign{}
\endhead
\bottomrule\noalign{}
\endlastfoot
\(a\) & \\
\(b\) & basic values \\
\(c\) & primitive values \\
\(d\) & declarations \\
\(e\) & expressions \\
\(f\) & builtin functions \\
\(g\) & \\
\(h\) & \\
\(i\) & (indexing) \\
\(j\) & (indexing) \\
\(k\) & \\
\(l\) & locations \\
\(m\) & modules \\
\(n\) & \\
\(o\) & \\
\(p\) & patterns \\
\(q\) & quantities \\
\(r\) & \\
\(s\) & \\
\(t\) & \\
\(u\) & \\
\(v\) & values \\
\(w\) & \\
\(x\) & variables \\
\(y\) & \\
\(z\) & \\
\end{longtable}

\begin{longtable}[]{@{}rl@{}}
\toprule\noalign{}
\endhead
\bottomrule\noalign{}
\endlastfoot
\(\alpha\) & type variables \\
\(\beta\) & basic types \\
\(\gamma\) & \\
\(\delta\) & \\
\(\varepsilon\) & borrow quantity \\
\(\zeta\) & \\
\(\eta\) & effect row \\
\(\theta\) & heap variables \\
\(\iota\) & \\
\(\kappa\) & kinds \\
\(\lambda\) & effect label \\
\(\mu\) & multiple quantity \\
\(\nu\) & owned quantity \\
\(\xi\) & \\
\(ο\) & \\
\(\pi\) & primitive types \\
\(\rho\) & row variable \\
\(\sigma\) & type schemas \\
\(\tau\) & types \\
\(\upsilon\) & \\
\(\varphi\) & \\
\(\chi\) & \\
\(\psi\) & \\
\(\omega\) & unrestricted quantity \\
\end{longtable}

We write:\\
\(\lbrack{\overline{\, x\:}}^{n}\text{ for an }\text{ordered}\text{ set of }x\text{es of length }n\text{, so }{x_{1},\ldots,x_{n}}.\rbrack\)

\section{Syntax}

\[\begin{array}{r}
\begin{array}{llll}
m & ::= & & \text{Modules}: \\
 & | & {{\overline{\, d\:}}^{\ast};e} & \text{– }\text{ main }
\end{array} \\
\begin{array}{llll}
d & ::= & & \text{Declarations}: \\
 & | & {\mathsf{\mathbf{\text{type}}}\ X = \langle{\overline{\, C^{n}\left( {\overline{\,\tau\:}}^{n} \right)\:}}^{m}\rangle} & \text{– }\text{ types}
\end{array} \\
\begin{array}{llll}
e & ::= & & \text{Expressions}: \\
 & | & x & \text{– }\text{ variable} \\
 & | & {{}^{{\overline{\, x\:}}^{\ast}}\left\{ e \right\}} & \text{– }\text{ borrow} \\
 & | & {|{\overline{\,{x\text{ :}^{q}\tau}\:}}^{n}|\ e_{0}} & \text{– }\text{ abstraction} \\
 & | & {e_{0}({\overline{\, e\:}}^{n})} & \text{– }\text{ application} \\
 & | & {(,{\overline{\, e\:}}^{n})} & \text{– }\text{ tuple} \\
 & | & {\mathsf{\mathbf{\text{val}}}^{q_{0}}\ {(,{\overline{\, x\:}}^{n})} = e_{0};\ e} & \text{– }\text{ split} \\
 & | & {C({\overline{\, e\:}}^{n})} & \text{– }\text{ variant} \\
 & | & {\mathsf{\mathbf{\text{match}}}^{q_{0}}\ e_{0}\ \{{\overline{\,{C({\overline{\, x\:}}^{n})}e\:}}^{m}\}} & \text{– }\text{ match}
\end{array} \\
\begin{array}{llll}
v & ::= & & \text{Values}: \\
 & | & {{{}^{{\overline{\, z\:}}^{\ast}}{|{\overline{\,{x\text{ :}^{q}\tau}\:}}^{n}|}}\ e_{0}} & \text{– }\text{ abstraction} \\
 & | & {(,{\overline{\, v\:}}^{n})} & \text{– }\text{ tuple} \\
 & | & {C({\overline{\, v\:}}^{n})} & \text{– }\text{ variant}
\end{array} \\
\begin{array}{llll}
b & ::= & & \text{Basic values}: \\
 & | & {(,{\overline{\, b\:}}^{n})} & \text{– }\text{ tuple} \\
 & | & {C({\overline{\, b\:}}^{n})} & \text{– }\text{ variant}
\end{array}
\end{array}\] Note we have \(b \subset v \subset e\)

\[\begin{array}{r}
\begin{array}{llll}
\tau & ::= & & \text{Types}: \\
 & | & \varphi & \text{– }\text{ arrow} \\
 & | & {(,{\overline{\,\tau\:}}^{n})} & \text{– }\text{ product} \\
 & | & {\langle,{\overline{\,{C({\overline{\,\tau\:}}^{n})}\:}}^{m}\rangle} & \text{– }\text{ sum} \\
 & | & {\text{List}({,\ }\tau)} & \text{– }\text{ list}
\end{array} \\
\begin{array}{llll}
\varphi & ::= & & \text{Function types}: \\
 & | & {({,\ }{\overline{\,{{}^{q}\tau}\:}}^{n})\tau_{0}} & \text{– }\text{ arrow}
\end{array} \\
\begin{array}{llll}
q & ::= & & \text{Quantities}: \\
 & | & \varepsilon & \text{– }\text{ borrowed} \\
 & | & 1 & \text{– }\text{ linear} \\
 & | & \omega & \text{– }\text{ unrestricted}
\end{array}
\end{array}\]

\subsection{Shorthands}

\[\begin{array}{llll}
\mu & \in & \left\{ \varepsilon,\omega \right\} & \text{– }\text{ multiple} \\
\nu & \in & \left\{ 1,\omega \right\} & \text{– }\text{ owned} \\
\pi & \in & \left\{ \varepsilon,1 \right\} & \text{– }\text{ parameter}
\end{array}\]

\subsection{Sugar}

\[\begin{array}{llll}
{\mathsf{\mathbf{\text{fun}}}\ x_{0}({\overline{\,{x\text{ :}^{q}\tau}\:}}^{n})\ e_{0};\ e} & ≔ & {\mathsf{\mathbf{\text{val}}}^{\omega}\ x_{0} = {|{\overline{\,{x\text{ :}^{q}\tau}\:}}^{n}|\ e_{0}};\ e} & \text{– }\text{ function declaration} \\
{\mathsf{\mathbf{\text{with}}}\ x_{0}\  \leftarrow {e_{0}({\overline{\, e\:}}^{n})};\ e} & ≔ & {e_{0}({\overline{\, e\:}}^{n},{|{x_{0}\text{ :}^{q}\tau}|\ e})} & \text{– }\text{ with binding }
\end{array}\]

About with-binding:
\(\lbrack\text{Should it be in an }\varepsilon\text{ context?}\rbrack\)

\section{Typing}

The algorithmic typing rules of \(\lambda^{\text{borrow}}\) are given in
between below text. The typing relation
\[\Gamma^{+}\  \vdash \ e^{+}\ :^{q^{+}}\ \tau^{-}\  \rightsquigarrow \ \Gamma^{-}\]
can be read as ``using expression \(e\) with quantity \(q\) can make use
of all the bindings in context \(\Gamma\), which yields type \(\tau\)
and a modified context \(\Gamma'\).''

\subsection{Variable lookup}

Variable lookup comes in two flavours. Linear bindings with quantity
\(1\) are looked up and removed from the context as shown in rule
\(\text{Var}_{1}\). Borrowed and unrestricted bindings with quantities
\(\varepsilon\) and \(\omega\) respectively, are looked up, but stay in
the resulting context. Rules \(\text{Var}_{\mu}\) define this for
\(\mu \in \left\{ \varepsilon,\omega \right\}\) simultaneously.
\[\begin{array}{r}
{\text{Var}_{1}\ \frac{\text{\quad\quad}\ }{\Gamma{,\ }{x\text{ :}^{1}\tau}\  \vdash \ x\ :^{1}\ \tau\  \rightsquigarrow \ \Gamma}\ } \\
{\text{Var}_{\mu}\ \frac{\text{\quad\quad}\ }{\Gamma{,\ }{x\text{ :}^{\mu}\tau}\  \vdash \ x\ :^{\mu}\ \tau\  \rightsquigarrow \ \Gamma{,\ }{x\text{ :}^{\mu}\tau}}\ \mu \in \left\{ \varepsilon,\omega \right\}}
\end{array}\]

We need a \emph{weakening} rule which states that unrestricted bindings
can be used linearly. Rules \(\text{Var}_{\pi}\) also state that
unrestricted bindings can be borrowed freely. Equivalently, we could
define weakening as a general rule on bindings instead of a rule for
variable lookup. However, this way our rule set would be
nondeterministic.
\[\text{Weak}\ \frac{\text{\quad\quad}{\Gamma\  \vdash \ x\ :^{\omega}\ \tau\  \rightsquigarrow \ \Gamma}}{\Gamma\  \vdash \ x\ :^{1}\ \tau\  \rightsquigarrow \ \Gamma}\ \]

\[\text{Var}_{\pi}\ \frac{\text{\quad\quad}\ }{\Gamma{,\ }{x\text{ :}^{\omega}\tau}\  \vdash \ x\ :^{\pi}\ \tau\  \rightsquigarrow \ \Gamma{,\ }{x\text{ :}^{\omega}\tau}}\ \pi \in \left\{ 1,\varepsilon \right\}\]

To allow linear bindings, that is bindings with quantity \(1\), to be
borrowed, we need to do additional bookkeeping. Borrows are only valid
in a lexical region. After this region ends, we restore the original
quantity on the binding.
\[\text{Borrow}_{1}\ \frac{\text{\quad\quad}{\Gamma_{0}{,\ }{\overline{\,{x\text{ :}^{\varepsilon}\tau}\:}}^{\ast}\  \vdash \ e_{0}\ :^{\left\lceil q \right\rceil}\ \tau_{0}\  \rightsquigarrow \ \Gamma_{1}{,\ }{\overline{\,{x\text{ :}^{\varepsilon}\tau}\:}}^{\ast}}}{\Gamma_{0}{,\ }{\overline{\,{x\text{ :}^{1}\tau}\:}}^{\ast}\  \vdash \ {{}^{{\overline{\, x\:}}^{\ast}}\left\{ e_{0} \right\}}\ :^{q}\ \tau_{0}\  \rightsquigarrow \ \Gamma_{1}{,\ }{\overline{\,{x\text{ :}^{1}\tau}\:}}^{\ast}}\ \tau \neq \varphi\]
Add explanation why linear functions cannot be borrowed. Here, we need
to take care borrowed bindings do not escape from this region.
Therefore, we \emph{lift} quantity \(q\) of the expression surroundings
to be an owned quantity. That is, borrowed expression surroundings are
lifted to unrestricted contexts, the two owning quantities stay the
same. The definition of lifting is as follows. \[\begin{array}{lll}
{\left\lceil \cdot \right\rceil:\text{ Quantity }\text{ Quantity}} & \left\lceil \varepsilon \right\rceil & = \\
\omega & \left\lceil q \right\rceil & = \\
q
\end{array}\]

Alternatively, we could enforce explicit borrowing of unrestricted
bindings, in the same way we do that for linear ones. We can change
\(\text{Borrow}_{1}\) to also include \(\omega\) bindings. Then, we need
to alter \(\text{Var}_{\pi}\) to remove \(\varepsilon\) as a free
borrow. \[\begin{array}{r}
{\text{Var}_{\text{weak}}\ \frac{\text{\quad\quad}\ }{\Gamma{,\ }{x\text{ :}^{\omega}\tau}\  \vdash \ x\ :^{1}\ \tau\  \rightsquigarrow \ \Gamma{,\ }{x\text{ :}^{\omega}\tau}}\ } \\
{\text{Borrow}_{\nu}\ \frac{\text{\quad\quad}{\Gamma_{0}{,\ }{\overline{\,{x\text{ :}^{\varepsilon}\tau}\:}}^{\ast}\  \vdash \ e_{0}\ :^{\left\lceil q \right\rceil}\ \tau_{0}\  \rightsquigarrow \ \Gamma_{1}{,\ }{\overline{\,{x\text{ :}^{\varepsilon}\tau}\:}}^{\ast}}}{\Gamma_{0}{,\ }{\overline{\,{x\text{ :}^{\nu}\tau}\:}}^{\ast}\  \vdash \ {{}^{{\overline{\, x\:}}^{\ast}}\left\{ e_{0} \right\}}\ :^{q}\ \tau_{0}\  \rightsquigarrow \ \Gamma_{1}{,\ }{\overline{\,{x\text{ :}^{\nu}\tau}\:}}^{\ast}}\ \nu \in \left\{ 1,\omega \right\} \land \tau \neq \varphi}
\end{array}\] This restricts the places where we can use explicit
borrowing. Is this good or bad? Could help in borrow inference: only
there were linear variables are used\ldots{}

\subsection{Functions}

For function abstraction, we have three cases, one for each quantity.
Depending on the quantity of the expression surroundings, anonymous
function blocks have access to different sets of bindings.

\begin{description}
\item[\(\varepsilon\)]
As we are in a borrowed expression surroundings, function blocks cannot
be returned nor stored: they are \emph{second-class}. Therefore, these
blocks have access to all borrowed bindings as well as all unrestricted
bindings. As they can be called multiple times (the code is borrowed and
can be used multiple times), we cannot allow usage of linear bindings.
\item[\(\omega\)]
For unrestricted expression surroundings the situation in different. As
these function blocks are owned, they \emph{can} be stored or returned.
Therefore, we need to make sure second-class bindings are not stored in
its closure. Borrowed bindings should not escape, only unrestricted
bindings are allowed.
\item[\(1\)]
Similarly, linear blocks are first-class and can be saved or returned,
so we cannot close over borrowed bindings. However, as we know that the
resulting closure can only be used \emph{once}, in this case we can also
allow access to linear bindings.
\end{description}

\[\begin{array}{r}
{\text{Abs}_{\varepsilon}\ \frac{\text{\quad\quad}{\Gamma_{0}^{\varepsilon}{,\ }\Gamma_{0}^{\omega}{,\ }{\overline{\,{x\text{ :}^{q}\tau}\:}}^{n}\  \vdash \ e_{0}\ :^{1}\ \tau_{0}\  \rightsquigarrow \ \Gamma_{1}}}{\Gamma_{0}\  \vdash \ {|{\overline{\,{x\text{ :}^{q}\tau}\:}}^{n}|\ e_{0}}\ :^{\varepsilon}\ {({,\ }{\overline{\,{{}^{q}\tau}\:}}^{n})\tau_{0}}\  \rightsquigarrow \ \Gamma_{0}^{1}{,\ }\Gamma_{1} \smallsetminus {\overline{\, x\:}}^{n}}\ } \\
{\text{Abs}_{1}\ \frac{\text{\quad\quad}{\Gamma_{0}^{1}{,\ }\Gamma_{0}^{\omega}{,\ }{\overline{\,{x\text{ :}^{q}\tau}\:}}^{n}\  \vdash \ e_{0}\ :^{1}\ \tau_{0}\  \rightsquigarrow \ \Gamma_{1}}}{\Gamma_{0}\  \vdash \ {|{\overline{\,{x\text{ :}^{q}\tau}\:}}^{n}|\ e_{0}}\ :^{1}\ {({,\ }{\overline{\,{{}^{q}\tau}\:}}^{n})\tau_{0}}\  \rightsquigarrow \ \Gamma_{0}^{\varepsilon}{,\ }\Gamma_{1} \smallsetminus {\overline{\, x\:}}^{n}}\ } \\
{\text{Abs}_{\omega}\ \frac{\text{\quad\quad}{\Gamma_{0}^{\omega}{,\ }{\overline{\,{x\text{ :}^{q}\tau}\:}}^{n}\  \vdash \ e_{0}\ :^{1}\ \tau_{0}\  \rightsquigarrow \ \Gamma_{1}}}{\Gamma_{0}\  \vdash \ {|{\overline{\,{x\text{ :}^{q}\tau}\:}}^{n}|\ e_{0}}\ :^{\omega}\ {({,\ }{\overline{\,{{}^{q}\tau}\:}}^{n})\tau_{0}}\  \rightsquigarrow \ \Gamma_{0}^{\varepsilon}{,\ }\Gamma_{0}^{1}{,\ }\Gamma_{1} \smallsetminus {\overline{\, x\:}}^{n}}\ }
\end{array}\]

For the bodies of the abstractions, there are two things to keep in
mind. First, we should not be allowed to return bindings that are
borrowed, so the quantity context of the body should be owning and
cannot be \(\varepsilon\). Second, the body \emph{should not know} how
many times its result will be used. Take for example the identity
function
\[{\mathsf{\mathbf{\text{fun}}}\ \text{identity}({x\text{ :}^{1}\tau})\ x;\ \ }\quad\]
which by definition {[}X{]} itself is unrestricted
\[{\mathsf{\mathbf{\text{val}}}^{\omega}\ \text{ identity} = {|{x\text{ :}^{1}\tau}|\ x};\ \ }.\]
This does not mean the body of the abstraction should be checked in an
\(\omega\) context. We have binding \(x\) linearly in the context, we
can return it, as it is owned. However, how the result will be used is a
responsibility of the caller. This means, it suffices to check the body
of an abstraction in a linear context!

To select bindings with the proper quantity from the context, we use
\emph{context filtering} which is defined as follows.
\[\begin{array}{lll}
{\Gamma^{q}:\text{ Context } \times \text{ Quantity }\text{ Context}} & \varnothing^{q} & = \\
\varnothing & \left( \Gamma{,\ }{x\text{ :}^{q}\tau} \right)^{q} & = \\
\Gamma^{q}{,\ }{x\text{ :}^{q}\tau} & \left( \Gamma{,\ }{x\text{ :}^{q'}\tau} \right)^{q} & = \\
\Gamma^{q}
\end{array}\]

When functions are applied in an expression surroundings of quantity
\(q\), the function itself needs to be available \(q\) times. Quantities
of the arguments are determined by the function's type signature.
\[\text{Wilt}\ \frac{{\Gamma_{0}\  \vdash \ e_{0}\ :^{1}\ {({,\ }{\overline{\,{{}^{q}\tau}\:}}^{n})\tau_{0}}\  \rightsquigarrow \ \Gamma_{1}}\text{\quad\quad}{\text{for each }i \in 1..n}\text{\quad\quad}{\Gamma_{i - 1} \Cap \Gamma_{i}\  \vdash \ e_{i}\ :^{q_{i}}\ \tau_{i}\  \rightsquigarrow \ \Gamma_{i + 1}}}{\Gamma_{0}\  \vdash \ \mathsf{\mathbf{\text{wilt}}}{e_{0}({\overline{\, e\:}}^{n})}\ :^{1}\ \tau_{0}\  \rightsquigarrow \ \Gamma_{n + 1}}\ \]
Is lowering really needed here? Calls can be just borrowed, that's
probably enough. However, do we restrict or allow something when calls
cannot be \(\omega\)?

\[\begin{array}{lll}
{\left\lfloor \cdot \right\rfloor:\text{ Quantity }\text{ Quantity}} & \left\lfloor \omega \right\rfloor & = \\
\varepsilon & \left\lfloor q \right\rfloor & = \\
q
\end{array}\]

Problems arise when linear arguments are passed to functions. Linear
arguments are owned and could be returned by a function. We need to make
sure the returned value can also only be used once.

The identity function simply returns its only parameter. The quantity of
this parameter cannot be \(\varepsilon\), as borrowed parameters cannot
be returned from functions, so it should be owned. We could pick
\(\omega\), but then we'd restrict ourselves because it cannot be used
on linear arguments. Therefore, the sensible choice is to pick \(1\).
\[\mathsf{\mathbf{\text{fun}}}\ \text{identity}({x\text{ :}^{1}\tau})\ x;\ \ \]

\(\begin{array}{r}
\mathsf{\mathbf{\text{val}}}^{1}\ a = 42;\  \\
\begin{array}{r}
\mathsf{\mathbf{\text{val}}}^{\omega}\ b = {\text{identity}(a)};\  \\
{(b,b)}
\end{array}
\end{array}\)

Variable \(a\) is used by passing it to \(\text{identity}\), the
returned value \(b\) however, is made available with quantity
\(\omega\). Therefore, we are allowed to create a pair of two \(b\)`s
and we've indirectly duplicated \(a\)\ldots{}

There are two solutions:

\begin{enumerate}
\item
  Functions with linear parameters can only be called in a linear
  quantity context.
\item
  Functions with linear parameters, when called in an unrestricted
  quantity context, need their arguments to be unrestricted as well.
\end{enumerate}

The first option would be a severe restriction, as many functions with
linear parameters should be usable on unrestricted arguments, and the
results would need to be used linearly. The second option seems more
logical.

\[\begin{array}{r}
{\text{App}_{\lambda}\ \frac{{\Gamma_{0}\  \vdash \ e_{0}\ :^{\varepsilon}\ {({,\ }{\overline{\,{{}^{q}\tau}\:}}^{n})\tau_{0}}\  \rightsquigarrow \ \Gamma_{1}}\text{\quad\quad}{\text{for each }i \in 1..n}\text{\quad\quad}{\Gamma_{i - 1} \Cap \Gamma_{i}\  \vdash \ e_{i}\ :^{q_{i}}\ \tau_{i}\  \rightsquigarrow \ \Gamma_{i + 1}}}{\Gamma_{0}\  \vdash \ {e_{0}({\overline{\, e\:}}^{n})}\ :^{\pi}\ \tau_{0}\  \rightsquigarrow \ \Gamma_{n + 1}}\ \pi \in \left\{ \varepsilon,1 \right\}} \\
{\text{App}_{\omega}\ \frac{{\Gamma_{0}\  \vdash \ e_{0}\ :^{\varepsilon}\ {({,\ }{\overline{\,{{}^{q}\tau}\:}}^{n})\tau_{0}}\  \rightsquigarrow \ \Gamma_{1}}\text{\quad\quad}{\text{for each }i \in 1..n}\text{\quad\quad}{\Gamma_{i - 1} \Cap \Gamma_{i}\  \vdash \ e_{i}\ :^{\left| q_{i} \right|}\ \tau_{i}\  \rightsquigarrow \ \Gamma_{i + 1}}}{\Gamma_{0}\  \vdash \ {e_{0}({\overline{\, e\:}}^{n})}\ :^{\omega}\ \tau_{0}\  \rightsquigarrow \ \Gamma_{n + 1}}\ }
\end{array}\]

\[\begin{array}{lll}
{| \cdot |:\text{ Quantity }\text{ Quantity}} & |1| & = \\
\omega & |q| & = \\
q
\end{array}\]

\subsection{Datatypes}

When creating datatypes, we need to store data so each subexpression in
constructors needs to be owned. Although we allow creating datatypes in
borrowed expression surroundings, we lift the surrounding quantity to
make sure stored data is owned.
\[\text{Construct}\ \frac{{\Delta\  \Vdash \ C\ :\ {({,\ }{\overline{\,\tau\:}}^{n})\tau_{0}}}\text{\quad\quad}{\text{for each }i \in 1..n}\text{\quad\quad}{\Gamma_{i - 1} \Cap \Gamma_{i}\  \vdash \ e_{i}\ :^{q}\ \tau_{i}\  \rightsquigarrow \ \Gamma_{i + 1}}}{\Gamma_{1}\  \vdash \ {C({\overline{\, e\:}}^{n})}\ :^{q}\ \tau_{0}\  \rightsquigarrow \ \Gamma_{n + 1}}\ \]
Note the similarities and differences between rules \(\text{Con}\) and
\(\text{App}\):
\(\lbrack\text{Both “lookup” the type of the function or constructor,
which directs the type of the arguments and the returntype of the application or construction.}\rbrack\)When
destructuring datatypes, we have two quantities to take into account:
\(\lbrack\text{The quantity in which the whole destructuring expression is going to be evaluated.
We call this the }\text{expression surrounding quantity}.\rbrack\)\[\text{Match}\ \frac{{\Gamma_{0}\  \vdash \ e_{0}\ :^{q_{0}}\ \tau_{0}\  \rightsquigarrow \ \Gamma'_{0}}\text{\quad\quad}{\text{for each }i \in 1..m}\text{\quad\quad}{\Delta\  \Vdash \ C_{i}\ :\ {({,\ }{\overline{\,\tau\:}}^{n_{i}})\tau_{0}}}\text{\quad\quad}{\Gamma'_{0}{,\ }{\overline{\,{x\text{ :}^{q_{0}}\tau}\:}}^{n_{i}}\  \vdash \ e_{i}\ :^{q}\ \tau\  \rightsquigarrow \ \Gamma_{i}}}{\Gamma_{0}\  \vdash \ {\mathsf{\mathbf{\text{match}}}^{q_{0}}\ e_{0}\ \{{\overline{\,{C({\overline{\, x\:}}^{n})}e\:}}^{m}\}}\ :^{q}\ \tau\  \rightsquigarrow \  \Cap_{i \in 1..m}\Gamma_{i}}\ \]

Because now we have multiple branches that can be taken, we need to
\emph{merge} the resulting contexts of each branch and remove the
freshly introduced bindings if still existing. As after branching,
contexts can only differ in removed linear bindings, merging two
contexts is simply set intersection. Check this! Aren't we passing
let-bound variables to the next argument?

\subsection{Tuples}

Tuples are \emph{unboxed}, which means they don't allocate space on the
heap and therefore do not store their components in the way datatypes
do. Therefore they their components can be borrowed.
\[\text{Pair}\ \frac{{\text{for each }i \in 1..n}\text{\quad\quad}{\Gamma_{i - 1} \Cap \Gamma_{i}\  \vdash \ e_{i}\ :^{q}\ \tau_{i}\  \rightsquigarrow \ \Gamma_{i + 1}}}{\Gamma_{1}\  \vdash \ {(,{\overline{\, e\:}}^{n})}\ :^{q}\ {(,{\overline{\,\tau\:}}^{n})}\  \rightsquigarrow \ \Gamma_{n + 1}}\ \]
Due to this design decision, we can define value binding in terms of
tuple creation and destruction. Otherwise we would not be able to
\emph{reborrow} a variable, that is, bind a borrowed variable to a new
name.

% \begin{Shaded}
% \begin{Highlighting}[]
% \KeywordTok{fun} \FunctionTok{reborrow}\OperatorTok{(}\NormalTok{ε }\VariableTok{x} \OperatorTok{:} \DataTypeTok{a}\OperatorTok{)}
%       \KeywordTok{val}\NormalTok{ ε }\FunctionTok{y} \OperatorTok{=}\NormalTok{ x}
%       \OperatorTok{...}

% \end{Highlighting}
% \end{Shaded}

\[\begin{array}{llll}
{\mathsf{\mathbf{\text{val}}}^{q}\ x = e;\ c} & ≔ & {\mathsf{\mathbf{\text{val}}}^{q}\ (x) = (e);\ c} & \text{– }\text{ value binding}
\end{array}\]

For the splitting of tuples, we ask an expression \(e_{0}\) to be
available for quantity \(q_{0}\). The resulting bindings \(x_{1}\) and
\(x_{2}\) are then made available for the same quantity \(q_{0}\) in the
remaining part of the program. Note we need to remove these bindings
from the resulting context \(\Gamma_{2}\) if they still exist. This is
similar to the way we handle matching on datatypes.
\[\text{Split}\ \frac{{\Gamma_{0}\  \vdash \ e_{0}\ :^{q_{0}}\ {(,{\overline{\,\tau\:}}^{n})}\  \rightsquigarrow \ \Gamma_{1}}\text{\quad\quad}{\Gamma_{1}{,\ }{\overline{\,{x\text{ :}^{q_{0}}\tau}\:}}^{n}\  \vdash \ e\ :^{q}\ \tau\  \rightsquigarrow \ \Gamma_{2}}}{\Gamma_{0}\  \vdash \ {\mathsf{\mathbf{\text{split}}}^{q_{0}}\ {(,{\overline{\, x\:}}^{n})} = e_{0};\ e}\ :^{q}\ \tau\  \rightsquigarrow \ \Gamma_{2} \smallsetminus {\overline{\, x\:}}^{n}}\ \]

\subsection{Built-ins}

Note that pre-defined functions, as well as any top-level functions, can
be used unrestricted. So their body is checked in a surrounding of
quantity \(\omega\). \[\begin{array}{r}
\begin{array}{llll}
\text{ fold}_{q} & : & {({{}^{q}{\text{List}({,\ }\tau_{1})}}{,\ }{{}^{1}\tau_{2}}{,\ }{{}^{\varepsilon}{({{}^{q}\tau_{1}}{,\ }{{}^{1}\tau_{2}})\tau_{2}}})\tau_{2}} & \text{– }\text{ fold list}
\end{array} \\
\begin{array}{llll}
\text{ Nil}_{\tau} & : & {\text{List}({,\ }\tau)} & \text{– }\text{ nil list} \\
\text{ Cons} & : & {(\tau{,\ }{\text{List}({,\ }\tau)}){\text{List}({,\ }\tau)}} & \text{– }\text{ cons list}
\end{array} \\
\begin{array}{llll}
\text{ Bool} & \text{:=} & {\langle\text{False}(),\text{ True}()\rangle} & \text{– }\text{ boolean type} \\
\text{ Option}(\tau) & \text{:=} & {\langle\text{None}(),\text{ Some}(\tau)\rangle} & \text{– }\text{ option type} \\
\text{ Result}\left( \tau_{1},\tau_{2} \right) & \text{:=} & {\langle\text{Wrong}\left( \tau_{1} \right),\text{ Right}\left( \tau_{2} \right)\rangle} & \text{– }\text{ either type}
\end{array}
\end{array}\]

\section{Overview}

\subsection{Variables}

\[\begin{array}{r}
{\text{Var}_{1}\ \frac{\text{\quad\quad}\ }{\Gamma{,\ }{x\text{ :}^{1}\tau}\  \vdash \ x\ :^{1}\ \tau\  \rightsquigarrow \ \Gamma}\ } \\
{\text{Var}_{\mu}\ \frac{\text{\quad\quad}\ }{\Gamma{,\ }{x\text{ :}^{\mu}\tau}\  \vdash \ x\ :^{\mu}\ \tau\  \rightsquigarrow \ \Gamma{,\ }{x\text{ :}^{\mu}\tau}}\ \mu \in \left\{ \varepsilon,\omega \right\}} \\
{\text{Var}_{\pi}\ \frac{\text{\quad\quad}\ }{\Gamma{,\ }{x\text{ :}^{\omega}\tau}\  \vdash \ x\ :^{\pi}\ \tau\  \rightsquigarrow \ \Gamma{,\ }{x\text{ :}^{\omega}\tau}}\ \pi \in \left\{ 1,\varepsilon \right\}} \\
 \\
{\text{Borrow}_{1}\ \frac{\text{\quad\quad}{\Gamma_{0}{,\ }{\overline{\,{x\text{ :}^{\varepsilon}\tau}\:}}^{\ast}\  \vdash \ e_{0}\ :^{\left\lceil q \right\rceil}\ \tau_{0}\  \rightsquigarrow \ \Gamma_{1}{,\ }{\overline{\,{x\text{ :}^{\varepsilon}\tau}\:}}^{\ast}}}{\Gamma_{0}{,\ }{\overline{\,{x\text{ :}^{1}\tau}\:}}^{\ast}\  \vdash \ {{}^{{\overline{\, x\:}}^{\ast}}\left\{ e_{0} \right\}}\ :^{q}\ \tau_{0}\  \rightsquigarrow \ \Gamma_{1}{,\ }{\overline{\,{x\text{ :}^{1}\tau}\:}}^{\ast}}\ \tau \neq \varphi}
\end{array}\]

\subsection{Uncurried}

\[\begin{array}{r}
{\text{Abs}_{\varepsilon}\ \frac{\text{\quad\quad}{\Gamma_{0}^{\varepsilon}{,\ }\Gamma_{0}^{\omega}{,\ }{\overline{\,{x\text{ :}^{q}\tau}\:}}^{n}\  \vdash \ e_{0}\ :^{1}\ \tau_{0}\  \rightsquigarrow \ \Gamma_{1}}}{\Gamma_{0}\  \vdash \ {|{\overline{\,{x\text{ :}^{q}\tau}\:}}^{n}|\ e_{0}}\ :^{\varepsilon}\ {({,\ }{\overline{\,{{}^{q}\tau}\:}}^{n})\tau_{0}}\  \rightsquigarrow \ \Gamma_{0}^{1}{,\ }\Gamma_{1} \smallsetminus {\overline{\, x\:}}^{n}}\ } \\
{\text{Abs}_{1}\ \frac{\text{\quad\quad}{\Gamma_{0}^{1}{,\ }\Gamma_{0}^{\omega}{,\ }{\overline{\,{x\text{ :}^{q}\tau}\:}}^{n}\  \vdash \ e_{0}\ :^{1}\ \tau_{0}\  \rightsquigarrow \ \Gamma_{1}}}{\Gamma_{0}\  \vdash \ {|{\overline{\,{x\text{ :}^{q}\tau}\:}}^{n}|\ e_{0}}\ :^{1}\ {({,\ }{\overline{\,{{}^{q}\tau}\:}}^{n})\tau_{0}}\  \rightsquigarrow \ \Gamma_{0}^{\varepsilon}{,\ }\Gamma_{1} \smallsetminus {\overline{\, x\:}}^{n}}\ } \\
{\text{Abs}_{\omega}\ \frac{\text{\quad\quad}{\Gamma_{0}^{\omega}{,\ }{\overline{\,{x\text{ :}^{q}\tau}\:}}^{n}\  \vdash \ e_{0}\ :^{1}\ \tau_{0}\  \rightsquigarrow \ \Gamma_{1}}}{\Gamma_{0}\  \vdash \ {|{\overline{\,{x\text{ :}^{q}\tau}\:}}^{n}|\ e_{0}}\ :^{\omega}\ {({,\ }{\overline{\,{{}^{q}\tau}\:}}^{n})\tau_{0}}\  \rightsquigarrow \ \Gamma_{0}^{\varepsilon}{,\ }\Gamma_{0}^{1}{,\ }\Gamma_{1} \smallsetminus {\overline{\, x\:}}^{n}}\ } \\
 \\
{\text{App}_{\lambda}\ \frac{{\Gamma_{0}\  \vdash \ e_{0}\ :^{\varepsilon}\ {({,\ }{\overline{\,{{}^{q}\tau}\:}}^{n})\tau_{0}}\  \rightsquigarrow \ \Gamma_{1}}\text{\quad\quad}{\text{for each }i \in 1..n}\text{\quad\quad}{\Gamma_{i - 1} \Cap \Gamma_{i}\  \vdash \ e_{i}\ :^{q_{i}}\ \tau_{i}\  \rightsquigarrow \ \Gamma_{i + 1}}}{\Gamma_{0}\  \vdash \ {e_{0}({\overline{\, e\:}}^{n})}\ :^{\pi}\ \tau_{0}\  \rightsquigarrow \ \Gamma_{n + 1}}\ \pi \in \left\{ \varepsilon,1 \right\}} \\
{\text{App}_{\omega}\ \frac{{\Gamma_{0}\  \vdash \ e_{0}\ :^{\varepsilon}\ {({,\ }{\overline{\,{{}^{q}\tau}\:}}^{n})\tau_{0}}\  \rightsquigarrow \ \Gamma_{1}}\text{\quad\quad}{\text{for each }i \in 1..n}\text{\quad\quad}{\Gamma_{i - 1} \Cap \Gamma_{i}\  \vdash \ e_{i}\ :^{\left| q_{i} \right|}\ \tau_{i}\  \rightsquigarrow \ \Gamma_{i + 1}}}{\Gamma_{0}\  \vdash \ {e_{0}({\overline{\, e\:}}^{n})}\ :^{\omega}\ \tau_{0}\  \rightsquigarrow \ \Gamma_{n + 1}}\ } \\
 \\
 \\
{\text{Pair}\ \frac{{\text{for each }i \in 1..n}\text{\quad\quad}{\Gamma_{i - 1} \Cap \Gamma_{i}\  \vdash \ e_{i}\ :^{q}\ \tau_{i}\  \rightsquigarrow \ \Gamma_{i + 1}}}{\Gamma_{1}\  \vdash \ {(,{\overline{\, e\:}}^{n})}\ :^{q}\ {(,{\overline{\,\tau\:}}^{n})}\  \rightsquigarrow \ \Gamma_{n + 1}}\ } \\
{\text{Split}\ \frac{{\Gamma_{0}\  \vdash \ e_{0}\ :^{q_{0}}\ {(,{\overline{\,\tau\:}}^{n})}\  \rightsquigarrow \ \Gamma_{1}}\text{\quad\quad}{\Gamma_{1}{,\ }{\overline{\,{x\text{ :}^{q_{0}}\tau}\:}}^{n}\  \vdash \ e\ :^{q}\ \tau\  \rightsquigarrow \ \Gamma_{2}}}{\Gamma_{0}\  \vdash \ {\mathsf{\mathbf{\text{split}}}^{q_{0}}\ {(,{\overline{\, x\:}}^{n})} = e_{0};\ e}\ :^{q}\ \tau\  \rightsquigarrow \ \Gamma_{2} \smallsetminus {\overline{\, x\:}}^{n}}\ } \\
 \\
{\text{Construct}\ \frac{{\Delta\  \Vdash \ C\ :\ {({,\ }{\overline{\,\tau\:}}^{n})\tau_{0}}}\text{\quad\quad}{\text{for each }i \in 1..n}\text{\quad\quad}{\Gamma_{i - 1} \Cap \Gamma_{i}\  \vdash \ e_{i}\ :^{q}\ \tau_{i}\  \rightsquigarrow \ \Gamma_{i + 1}}}{\Gamma_{1}\  \vdash \ {C({\overline{\, e\:}}^{n})}\ :^{q}\ \tau_{0}\  \rightsquigarrow \ \Gamma_{n + 1}}\ } \\
{\text{Match}\ \frac{{\Gamma_{0}\  \vdash \ e_{0}\ :^{q_{0}}\ \tau_{0}\  \rightsquigarrow \ \Gamma'_{0}}\text{\quad\quad}{\text{for each }i \in 1..m}\text{\quad\quad}{\Delta\  \Vdash \ C_{i}\ :\ {({,\ }{\overline{\,\tau\:}}^{n_{i}})\tau_{0}}}\text{\quad\quad}{\Gamma'_{0}{,\ }{\overline{\,{x\text{ :}^{q_{0}}\tau}\:}}^{n_{i}}\  \vdash \ e_{i}\ :^{q}\ \tau\  \rightsquigarrow \ \Gamma_{i}}}{\Gamma_{0}\  \vdash \ {\mathsf{\mathbf{\text{match}}}^{q_{0}}\ e_{0}\ \{{\overline{\,{C({\overline{\, x\:}}^{n})}e\:}}^{m}\}}\ :^{q}\ \tau\  \rightsquigarrow \  \Cap_{i \in 1..m}\Gamma_{i}}\ }
\end{array}\]

\subsection{Curried}

\[\begin{array}{r}
{\text{Abs}_{\varepsilon}\ \frac{\text{\quad\quad}{\Gamma_{0}^{\varepsilon}{,\ }\Gamma_{0}^{\omega}{,\ }{x_{1}\text{ :}^{q_{1}}\tau_{1}}\  \vdash \ e_{0}\ :^{1}\ \tau_{0}\  \rightsquigarrow \ \Gamma_{1}}}{\Gamma_{0}\  \vdash \ {|{x_{1}\text{ :}^{q_{1}}\tau_{1}}|\ e_{0}}\ :^{\varepsilon}\ {({,\ }{{}^{q_{1}}\tau_{1}})\tau_{0}}\  \rightsquigarrow \ \Gamma_{0}^{1}{,\ }\Gamma_{1} \smallsetminus x_{1}}\ } \\
{\text{Abs}_{1}\ \frac{\text{\quad\quad}{\Gamma_{0}^{1}{,\ }\Gamma_{0}^{\omega}{,\ }{x_{1}\text{ :}^{q_{1}}\tau_{1}}\  \vdash \ e_{0}\ :^{1}\ \tau_{0}\  \rightsquigarrow \ \Gamma_{1}}}{\Gamma_{0}\  \vdash \ {|{x_{1}\text{ :}^{q_{1}}\tau_{1}}|\ e_{0}}\ :^{1}\ {({,\ }{{}^{q_{1}}\tau_{1}})\tau_{0}}\  \rightsquigarrow \ \Gamma_{0}^{\varepsilon}{,\ }\Gamma_{1} \smallsetminus x_{1}}\ } \\
{\text{Abs}_{\omega}\ \frac{\text{\quad\quad}{\Gamma_{0}^{\omega}{,\ }{x_{1}\text{ :}^{q_{1}}\tau_{1}}\  \vdash \ e_{0}\ :^{1}\ \tau_{0}\  \rightsquigarrow \ \Gamma_{1}}}{\Gamma_{0}\  \vdash \ {|{x_{1}\text{ :}^{q_{1}}\tau_{1}}|\ e_{0}}\ :^{\omega}\ {({,\ }{{}^{q_{1}}\tau_{1}})\tau_{0}}\  \rightsquigarrow \ \Gamma_{0}^{\varepsilon}{,\ }\Gamma_{0}^{1}{,\ }\Gamma_{1} \smallsetminus x_{1}}\ } \\
 \\
{\text{App}_{\lambda}\ \frac{{\Gamma_{0}\  \vdash \ e_{0}\ :^{\varepsilon}\ {({,\ }{{}^{q_{1}}\tau_{1}})\tau_{0}}\  \rightsquigarrow \ \Gamma_{1}}\text{\quad\quad}{\Gamma_{1}\  \vdash \ e_{1}\ :^{q_{1}}\ \tau_{1}\  \rightsquigarrow \ \Gamma_{2}}}{\Gamma_{0}\  \vdash \ {e_{0}(e_{1})}\ :^{\pi}\ \tau_{0}\  \rightsquigarrow \ \Gamma_{2}}\ \pi \in \left\{ \varepsilon,1 \right\}} \\
{\text{App}_{\omega}\ \frac{{\Gamma_{0}\  \vdash \ e_{0}\ :^{\varepsilon}\ {({,\ }{{}^{q_{1}}\tau_{1}})\tau_{0}}\  \rightsquigarrow \ \Gamma_{1}}\text{\quad\quad}{\Gamma_{1}\  \vdash \ e_{1}\ :^{\left| q_{1} \right|}\ \tau_{1}\  \rightsquigarrow \ \Gamma_{2}}}{\Gamma_{0}\  \vdash \ {e_{0}(e_{1})}\ :^{\omega}\ \tau_{0}\  \rightsquigarrow \ \Gamma_{2}}\ } \\
 \\
 \\
{\text{Pair}\ \frac{{\Gamma_{1}\  \vdash \ e_{1}\ :^{q}\ \tau_{1}\  \rightsquigarrow \ \Gamma_{2}}\text{\quad\quad}{\Gamma_{2}\  \vdash \ e_{2}\ :^{q}\ \tau_{2}\  \rightsquigarrow \ \Gamma_{3}}}{\Gamma_{1}\  \vdash \ {(e_{1},e_{2})}\ :^{q}\ {(\tau_{1},\tau_{2})}\  \rightsquigarrow \ \Gamma_{3}}\ } \\
{\text{Split}\ \frac{{\Gamma_{0}\  \vdash \ e_{0}\ :^{q_{0}}\ {(\tau_{1},\tau_{2})}\  \rightsquigarrow \ \Gamma_{1}}\text{\quad\quad}{\Gamma_{1}{,\ }{x_{1}\text{ :}^{q_{0}}\tau_{1}}{,\ }{x_{2}\text{ :}^{q_{0}}\tau_{2}}\  \vdash \ e\ :^{q}\ \tau\  \rightsquigarrow \ \Gamma_{2}}}{\Gamma_{0}\  \vdash \ {\mathsf{\mathbf{\text{split}}}^{q_{0}}\ {(x_{1},x_{2})} = e_{0};\ e}\ :^{q}\ \tau\  \rightsquigarrow \ \Gamma_{2} \smallsetminus x_{1} \smallsetminus x_{2}}\ } \\
 \\
{\text{Construct}\ \frac{{\Delta\  \Vdash \ C\ :\ {({,\ }\tau_{1})\tau_{0}}}\text{\quad\quad}{\Gamma_{1}\  \vdash \ e_{1}\ :^{q}\ \tau_{1}\  \rightsquigarrow \ \Gamma_{2}}}{\Gamma_{1}\  \vdash \ {C(e_{1})}\ :^{q}\ \tau_{0}\  \rightsquigarrow \ \Gamma_{2}}\ } \\
{\text{Match}\ \frac{{\Gamma_{0}\  \vdash \ e_{0}\ :^{q_{0}}\ \tau_{0}\  \rightsquigarrow \ \Gamma'_{0}}\text{\quad\quad}{\text{for each }i \in 1..2}\text{\quad\quad}{\Delta\  \Vdash \ C_{i}\ :\ {({,\ }\tau_{i})\tau_{0}}}\text{\quad\quad}{\Gamma'_{0}{,\ }{x_{i}\text{ :}^{q_{0}}\tau_{i}}\  \vdash \ e_{i}\ :^{q}\ \tau\  \rightsquigarrow \ \Gamma_{i}}}{\Gamma_{0}\  \vdash \ {\mathsf{\mathbf{\text{match}}}^{q_{0}}\ e_{0}\ \{{\overline{\,{C(x)}e\:}}^{2}\}}\ :^{q}\ \tau\  \rightsquigarrow \ \Gamma_{1} \Cap \Gamma_{2}}\ }
\end{array}\]
